\documentclass[a4paper,11pt,twoside]{article}
\usepackage[utf8]{inputenc}
\usepackage[english,greek]{babel}
\usepackage{alphabeta,amsmath}
\usepackage{nimbusserif}
\usepackage[T1]{fontenc}
\usepackage[outer=1.50cm, inner=2.00cm, top=2.00cm, bottom=2.00cm]{geometry}
\usepackage{enumitem,tcolorbox,xcolor}

\newcommand{\kerkissans}[1]{{\fontfamily{maksf}\selectfont {#1}}}

\title{\vspace{-2cm}\kerkissans{\textbf{Επαναληπτικά Θέματα Γ Λυκείου - Θεωρία (Θέμα Α)}}\vspace{-0.5cm}}
\date{}

\begin{document}
\maketitle

\noindent \kerkissans{\textbf{ΕΡΩΤΗΣΕΙΣ ΣΩΣΤΟ - ΛΑΘΟΣ}}

\vspace{0.3cm}
\noindent Χαρακτηρίστε τις παρακάτω προτάσεις γράφοντας δίπλα στο γράμμα που αντιστοιχεί σε κάθε πρόταση τη λέξη \textbf{Σωστό}, αν η πρόταση είναι σωστή, ή \textbf{Λάθος}, αν η πρόταση είναι λανθασμένη. Όλες οι ερωτήσεις ανταποκρίνονται στην ύλη των Πανελλαδικών Εξετάσεων και κάποιες κρύβουν συχνές "παγίδες".

\begin{enumerate}[label=\textbf{\arabic*.}]
    \item Αν $f'(x) = 0$ για κάθε $x$ στο πεδίο ορισμού της, τότε η $f$ είναι υποχρεωτικά σταθερή συνάρτηση σε όλο το πεδίο ορισμού της.
    \item Αν μια συνάρτηση $f$ είναι συνεχής στο $[\alpha, \beta]$ και ισχύει $f(\alpha) \cdot f(\beta) > 0$, τότε η $f$ σίγουρα δεν έχει ρίζα στο διάστημα $(\alpha, \beta)$.
    \item Κάθε συνεχής συνάρτηση με πεδίο ορισμού ένα ανοικτό διάστημα $(\alpha, \beta)$ έχει πάντοτε μέγιστη και ελάχιστη τιμή.
    \item Αν μια συνάρτηση $f$ είναι γνησίως αύξουσα σε ένα διάστημα $\Delta$, τότε υποχρεωτικά ισχύει $f'(x) > 0$ για κάθε εσωτερικό σημείο $x$ του $\Delta$, στο οποίο είναι παραγωγίσιμη.
    \item Αν διευκρινίζεται ότι $\lim\limits_{x \to x_0} f(x) = +\infty$, τότε ισχύει απαραίτητα ότι $f(x) > 0$ για κάθε $x$ πολύ κοντά στο $x_0$.
    \item Εάν η $f$ παρουσιάζει τοπικό ακρότατο στο εσωτερικό σημείο $x_0$ και είναι παραγωγίσιμη σε αυτό, τότε η εφαπτομένη της γραφικής της παράστασης στο $x_0$ είναι οριζόντια.
    \item Αν η συνάρτηση $f$ είναι $1-1$ στο πεδίο ορισμού της, τότε είναι υποχρεωτικά και γνησίως μονότονη.
    \item Έστω μια συνεχής συνάρτηση $f$ ορισμένη στο $[\alpha, \beta]$. Αν ισχύει $\int_\alpha^\beta f(x) dx = 0$ και $f(x) \ge 0$ για κάθε $x \in [\alpha, \beta]$, τότε αναγκαία ισχύει $f(x) = 0$ για κάθε $x \in [\alpha, \beta]$.
    \item Η ιδιότητα του ολοκληρώματος $\int_\alpha^\beta f(x) dx = \int_\alpha^\gamma f(x) dx + \int_\gamma^\beta f(x) dx$ ισχύει μόνο στην περίπτωση που το $\gamma$ ανήκει γνήσια μεταξύ του $\alpha$ και του $\beta$ ($\alpha < \gamma < \beta$).
    \item Αν ισχύει αυστηρά $f(x) < g(x)$ περιοχικά γύρω από το $x_0$, τότε υποχρεωτικά ισχύει και $\lim\limits_{x \to x_0} f(x) < \lim\limits_{x \to x_0} g(x)$ (εφόσον τα όρια υπάρχουν κανονικά στο $\mathbb{R}$).
    \item Κάθε συνάρτηση που είναι συνεχής σε ένα σημείο $x_0$, είναι απαραίτητα και παραγωγίσιμη στο σημείο αυτό.
    \item Αν $f''(x) > 0$ για κάθε $x \in \mathbb{R}$, τότε η συνάρτηση $f$ παρουσιάζει υποχρεωτικά ολικό ελάχιστο.
    \item Το πεδίο ορισμού της σύνθεσης $f \circ g$ ταυτίζεται πάντοτε με την τομή των πεδίων ορισμού της $f$ και της $g$.
    \item Κάθε κατακόρυφη ευθεία του επιπέδου έχει το πολύ ένα κοινό σημείο με τη γραφική παράσταση οποιασδήποτε συνάρτησης $f$.
    \item Για τη συνάρτηση $f(x) = |x|$, η παράγωγος δίνεται από τον τύπο $f'(x) = \dfrac{x}{|x|}$ για κάθε $x \ne 0$.
    \item Αν η $f$ είναι συνεχής στο $[\alpha, \beta]$ και ισχύει $\int_\alpha^\beta f(x) dx > 0$, τότε υποχρεωτικά επαληθεύεται ότι $f(x) > 0$ για κάθε $x \in [\alpha, \beta]$.
    \item Αν μια συνάρτηση είναι άρτια και παραγωγίσιμη σε όλο το $\mathbb{R}$, τότε η πρώτη της παράγωγος είναι συνάρτηση περιττή.
    \item Είναι πιθανό η γραφική παράσταση μιας συνάρτησης να τέμνει την κατακόρυφη ασύμπτωτή της.
    \item Είναι απολύτως δυνατό η γραφική παράσταση μιας συνάρτησης να τέμνει την πλάγια ή την οριζόντια ασύμπτωτή της.
    \item Για οποιαδήποτε συνάρτηση, αν ισχύει $\lim\limits_{x \to x_0} |f(x)| = 0$, τότε συνεπάγεται ότι $\lim\limits_{x \to x_0} f(x) = 0$.
    \item Αν ισχύει $\lim\limits_{x \to x_0} |f(x)| = \ell$ (όπου $\ell > 0$), τότε υποχρεωτικά θα ισχύει $\lim\limits_{x \to x_0} f(x) = \ell$ ή $\lim\limits_{x \to x_0} f(x) = -\ell$.
    \item Κάθε πολυωνυμική συνάρτηση τρίτου βαθμού (π.χ. $f(x) = ax^3+bx^2+cx+d, a \ne 0$) έχει πάντοτε ακριβώς ένα σημείο καμπής.
    \item Έστω συνάρτηση $f$ συνεχής στο $x_0$. Αν $f'(x) > 0$ για κάθε $x \in (\alpha, x_0) \cup (x_0, \beta)$, τότε η $f$ είναι γνησίως αύξουσα σε ολόκληρο το διάστημα $(\alpha, \beta)$.
    \item Αν $f'(x) = g'(x)$ για κάθε $x$ στο κοινό πεδίο ορισμού τους $\Delta$, τότε υποχρεωτικά $f(x) = g(x) + c$ για όλα τα $x \in \Delta$ (γενικά για οποιοδήποτε σύνολο \Delta).
    \item Οι γραφικές παραστάσεις $C_f$ και $C_{f^{-1}}$ δύο αντίστροφων συναρτήσεων είναι συμμετρικές ως προς την ευθεία διχοτόμο $y=x$.
    \item Το εμβαδόν του χωρίου μεταξύ των καμπυλών δύο συνεχών συναρτήσεων $f, g$ δίνεται πάντοτε από τον τύπο $E = \int_\alpha^\beta \big(f(x) - g(x)\big) dx$.
    \item Το άθροισμα $f+g$ δύο γνησίως αυξουσών συναρτήσεων είναι κι αυτό γνησίως αύξουσα συνάρτηση (στο κοινό πεδίο ορισμού).
    \item Η συνάρτηση που ορίζεται μέσω του ολοκληρώματος $F(x) = \int_a^x f(t)dt$ είναι υποχρεωτικά σταθερή, μόνο αν η $f$ είναι η μηδενική συνάρτηση.
    \item Αν μια συνάρτηση $f$ είναι συνεχής στο $x_0=0$ και γνωρίζουμε ότι υπάρχει στο $\mathbb{R}$ το όριο $\lim\limits_{x \to 0} \dfrac{f(x)}{x}$, τότε είναι αναγκαστικό γνώρισμα ότι $f(0)=0$.
    \item Γράφοντας $\int_\alpha^\beta f(t) dt$ ή $\int_\alpha^\beta f(x) dx$ αναφερόμαστε ακριβώς στην ίδια πραγματική σταθερά (ίδιος αριθμός).
    \item Εάν μια συνεχής συνάρτηση $f$ δεν μηδενίζεται σε κανένα σημείο ενός κλειστού διαστήματος $[\alpha, \beta]$, τότε σίγουρα διατηρεί σταθερό πρόσημο σε αυτό.
    \item Υπάρχει συνάρτηση $f$, παραγωγίσιμη σε όλο το $\mathbb{R}$, της οποίας η παράγωγος να περιέχει απόλυτη τιμή της μορφής $f'(x) = |x|$.
    \item Εάν η εφαπτομένη σε ένα σημείο $M$ της γραφικής παράστασης διαπερνά την καμπύλη, τότε το $M$ ονομάζεται σημείο καμπής.
    \item Ισχύει γενικά ότι αν $\lim\limits_{x \to x_0} \big(f(x)g(x)\big) = 0$, τότε υποχρεωτικά $\lim\limits_{x \to x_0} f(x) = 0$ ή $\lim\limits_{x \to x_0} g(x) = 0$.
    \item Κάθε αντιστρέψιμη ($1-1$) συνάρτηση, της οποίας το πεδίο ορισμού είναι διάστημα, είναι απαραίτητα και συνεχής.
    \item Αν στο σημείο $x_0$ έχουμε $f'(x_0) = 0$ και $f''(x_0) = 0$, τότε το $x_0$ σίγουρα αποκλείεται να αποτελεί τοπικό ακρότατο της συνάρτησης.
    \item Αν η γραφική παράσταση μιας συνάρτησης έχει οριζόντια ασύμπτωτη στο $+\infty$, τότε αποκλείεται η τιμή της $f(x)$ να τείνει στο $+\infty$ όταν το $x \to +\infty$.
    \item Είναι εφικτό ένα τοπικό μέγιστο μιας συνάρτησης να λαμβάνει μικρότερη αριθμητική τιμή από ένα τοπικό ελάχιστο της ίδιας συνάρτησης.
    \item Κάθε συνεχής συνάρτηση $f: \mathbb{R} \to \mathbb{R}$, θετική στο $+\infty$ και αρνητική στο $-\infty$, περιλαμβάνει ρίζα.
    \item Το βασικό τριγωνομετρικό όριο για άπειρο είναι $\lim\limits_{x \to +\infty} \dfrac{\sin x}{x} = 1$.
\end{enumerate}

\newpage
\noindent \kerkissans{\Large{\textbf{ΑΠΑΝΤΗΣΕΙΣ \& ΔΙΚΑΙΟΛΟΓΗΣΗ (Για τον Διδάσκοντα)}}}

\vspace{0.5cm}
\begin{tcolorbox}[colback=black!5!white, colframe=black, title=\textbf{Κλειδί Απαντήσεων}]
\begin{multicols}{4}
\begin{itemize}[label=$-$]
    \item \textbf{1.} Λάθος
    \item \textbf{2.} Λάθος
    \item \textbf{3.} Λάθος
    \item \textbf{4.} Λάθος
    \item \textbf{5.} Σωστό
    \item \textbf{6.} Σωστό
    \item \textbf{7.} Λάθος
    \item \textbf{8.} Σωστό
    \item \textbf{9.} Λάθος
    \item \textbf{10.} Λάθος
    \item \textbf{11.} Λάθος
    \item \textbf{12.} Λάθος
    \item \textbf{13.} Λάθος
    \item \textbf{14.} Σωστό
    \item \textbf{15.} Σωστό
    \item \textbf{16.} Λάθος
    \item \textbf{17.} Σωστό
    \item \textbf{18.} Λάθος
    \item \textbf{19.} Σωστό
    \item \textbf{20.} Σωστό
    \item \textbf{21.} Λάθος
    \item \textbf{22.} Σωστό
    \item \textbf{23.} Σωστό
    \item \textbf{24.} Λάθος
    \item \textbf{25.} Σωστό
    \item \textbf{26.} Λάθος
    \item \textbf{27.} Σωστό
    \item \textbf{28.} Σωστό
    \item \textbf{29.} Σωστό
    \item \textbf{30.} Σωστό
    \item \textbf{31.} Σωστό
    \item \textbf{32.} Σωστό
    \item \textbf{33.} Σωστό
    \item \textbf{34.} Λάθος
    \item \textbf{35.} Λάθος
    \item \textbf{36.} Λάθος
    \item \textbf{37.} Σωστό
    \item \textbf{38.} Σωστό
    \item \textbf{39.} Σωστό
    \item \textbf{40.} Λάθος
\end{itemize}
\end{multicols}
\end{tcolorbox}

\vspace{0.3cm}
\noindent \textbf{Σημαντικές Δικαιολογήσεις (Παγίδες):}
\begin{itemize}
    \item \textbf{1:} Λάθος, διότι αν το πεδίο ορισμού είναι ένωση ξένων διαστημάτων (π.χ. $1/x$), η συνάρτηση μπορεί να έχει διαφορετικές σταθερές σε κάθε διάστημα. Πρέπει να είναι σε \textit{ένα μόνο διάστημα}.
    \item \textbf{2:} Το θεώρημα Bolzano εξασφαλίζει ρίζα όταν $f(a)f(b)<0$.  Όταν $f(a)f(b)>0$, τίποτα δεν απαγορεύει στη συνάρτηση να έχει ρίζες (π.χ. άρτιο πλήθος ριζών όπως η παραβολή).
    \item \textbf{4:} Μπορεί να μηδενίζεται σε συγκερκιμένα σημεία, π.χ. η $f(x)=x^3$ είναι γνησίως αύξουσα παρ'όλο που $x=0 \implies f'(0)=0$.
    \item \textbf{7:} Μια συνάρτηση $1-1$ μπορεί να διαιρείται σε κλάδους και να μην είναι μονότονη. (π.χ. η συνάρτηση $1/x$).
    \item \textbf{9:} Η ιδιότητα πρόσθεσης ολοκληρωμάτων ισχύει για ΟΛΑ τα $\alpha, \beta, \gamma \in \Delta$, ανεξαρτήτως διάταξης, αρκεί να ανήκουν στο πεδίο ορισμού της συνεχούς $f$. 
    \item \textbf{10:} Όταν $f(x) < g(x)$, είναι πιθανό στο όριο οι συναρτήσεις να "αγγίξουν" η μία την άλλη ($=$), π.χ. $\frac{1}{x} > 0$ αλλά το όριο στο $+\infty$ είναι $0=0$.
    \item \textbf{12:} Η συνάρτηση $e^x$ έχει συνεχώς $f''(x)>0$, αλλά δεν έχει κανένα ακρότατο.
    \item \textbf{21:} Το όριο μπορεί απλώς να μην υπάρχει (η συνάρτηση ταλαντώνεται μεταξύ $\ell$ και $-\ell$).
    \item \textbf{26:} Χρειάζεται το ολοκλήρωμα της απόλυτης τιμής $|f(x)-g(x)|$.
    \item \textbf{29:} Πολύ σωστό: αν $\lim \frac{f(x)}{x} \in \mathbb{R}$, τότε $\lim f(x) = \lim (x \cdot \frac{f(x)}{x}) = 0$. Αφού είναι συνεχής στο 0, τότε αναγκαστικά $f(0)=0$.
    \item \textbf{36:} Η $f(x)=x^4$ έχει $f'(0)=0$ και $f''(0)=0$, ωστόσο το $0$ είναι ολικό ελάχιστο.
    \item \textbf{38:} Αν η συνάρτηση δεν είναι σε συνεχόμενο πεδίο ή δίκλαδη, ή αν υπάρχουν πολλαπλά ακρότατα (π.χ $f(x) = x + \sin x$ - όχι αυτό δεν έχει), ένα τοπικό μέγιστο σε μια 'κορυφή' μπορεί να είναι κάτω από ένα 'βάθος' αλλού. 
    \item \textbf{40:} Το $\lim_{x\to +\infty} \frac{\sin x}{x}$ αποδεικνύεται 0 (Κριτήριο Παρεμβολής). Μόνο στο $x \to 0$ κάνει 1.
\end{itemize}

\end{document}
